\documentclass[11pt]{article}
\usepackage[letterpaper,margin=1.02in]{geometry}
\usepackage[T1]{fontenc}
\usepackage{lmodern}
\usepackage{amsmath,amssymb,amsthm,booktabs,microtype}
\usepackage[colorlinks=true,linkcolor=blue!45!black,urlcolor=blue!45!black,citecolor=blue!45!black]{hyperref}
\usepackage{xcolor}
\newtheorem{theorem}{Theorem}[section]
\newtheorem{lemma}[theorem]{Lemma}
\newtheorem{proposition}[theorem]{Proposition}
\newtheorem{corollary}[theorem]{Corollary}
\theoremstyle{definition}
\newtheorem{definition}[theorem]{Definition}
\newcommand{\adim}{\operatorname{adim}}
\newcommand{\HH}{H_2}
\newcommand{\B}{B}
\newcommand{\Ssum}{S}
\numberwithin{equation}{section}
\setlength{\emergencystretch}{2em}
\raggedbottom
\hypersetup{pdftitle={Exact Autonomy Dimension of Tetrahedral Windmills},
pdfsubject={An exact solution of the windmill question from the autonomy-dimension research note},
pdfkeywords={graph theory, random walk, lumpability, windmill, partition packing}}
\title{Exact Autonomy Dimension\\of Tetrahedral Windmills}
\author{}
\date{October 26, 2026}
\begin{document}
\maketitle

\begin{abstract}
Let $W_r$ be the graph formed from $r$ copies of $K_4$ sharing a common
vertex. Its autonomy dimension is the minimum number of binary
observations that separate all vertices and individually project simple
random walk to a Markov chain for every starting state. An earlier note
asked for the exact value, having established only a first-order
asymptotic formula. We answer that question for every $r$. The answer is
the inverse of an explicit binomial-sum capacity formula, with three cases
according to the coordinate count modulo three. The proof converts
autonomous codes into packings of three-block partitions, bounds these
packings by weights on subset sizes, and attains the bounds using a
classical almost-regular hypergraph decomposition theorem. We also obtain
the second-order formula
\[
 \adim(W_r)=
 \frac{\log_2(3r)+\tfrac12\log_2\log_2(3r)}{H_2(1/3)}+O(1).
\]
Integer certificates and explicitly constructed codes accompany the proof.
\end{abstract}

\noindent\textbf{Keywords.} Strong lumpability; autonomous coordinates;
windmill graph; set partitions; hypergraph decomposition.

\section{The question and its exact answer}

Let $G$ be a finite connected simple graph with at least two vertices.
A nonconstant function $f:V(G)\to\{0,1\}$ is \emph{autonomous} if the
fraction of neighbors of opposite color is a constant $a$ on $f^{-1}(0)$
and a constant $b$ on $f^{-1}(1)$. Its observed random walk has transition
matrix $\left(\begin{smallmatrix}1-a&a\\b&1-b\end{smallmatrix}\right)$ for
every starting state. The \emph{autonomy dimension} $\adim(G)$ is the
minimum number of autonomous binary functions whose joint values
distinguish all vertices; it is infinity if no such family exists.

Write $W_r$ for $r$ copies of $K_4$ with one vertex identified to a common
hub $o$. There are $3r+1$ vertices. Removing the hub leaves $r$ triangles.
The October 7 research note~\cite{previous}, Section~7, asked for the exact
autonomy dimension of this family. This is the specific open question
addressed here; it is a question from that research program, not a claim
to settle an independently established named conjecture.

For a fixed integer $d\geq2$, set
\[
 B_j=\binom dj,\qquad S_j=\sum_{i=0}^j\binom di,
\]
where $B_j=0$ and $S_j=0$ when $j<0$. Their dependence on $d$ is implicit.

\begin{theorem}[Exact capacity and dimension]\label{thm:main}
Let $M_d$ be the largest $r$ for which $W_r$ has a separating autonomous
code using at most $d$ coordinates. Then
\begin{equation}\label{eq:capacity}
 M_d=
 \begin{cases}
 S_{k-2}+\left\lfloor(2B_{k-1}+B_k)/3\right\rfloor,
     &d=3k,\quad k\geq1,\\[2pt]
 S_{k-1}+\left\lfloor B_k/2\right\rfloor,
     &d=3k+1,\quad k\geq1,\\[2pt]
 S_{k-2}+B_k+\left\lfloor3B_{k-1}/4\right\rfloor,
     &d=3k+2,\quad k\geq0.
 \end{cases}
\end{equation}
Consequently, for every integer $r\geq1$,
\begin{equation}\label{eq:dimension}
 \boxed{\adim(W_r)=\min\{d\geq2:r\leq M_d\}.}
\end{equation}
\end{theorem}

\begin{table}[ht]
\centering
\begin{tabular}{c|rrrrrrrrrrrr}
\toprule
$d$&2&3&4&5&6&7&8&9&10&11&12&13\\
\midrule
$M_d$&1&1&3&5&10&18&35&62&116&218&390&735\\
\bottomrule
\end{tabular}
\caption{Exact capacities. For example, $\adim(W_{18})=7$ and
$\adim(W_{19})=8$. The repeated capacity at $d=3$ means that no windmill
has autonomy dimension three.}
\end{table}

Formula~\eqref{eq:dimension} is an explicit evaluation rule involving only
binomial coefficients and integer arithmetic. It requires no search
over graph colorings. The proof below relies on Bryant's established
hypergraph balancing theorem~\cite{Bryant}; the decomposition theorem
itself is not new. Historical priority of the exact partition-packing
formula or the graph-theoretic invariant has not been established by a
comprehensive literature review.

\section{From graph coordinates to partitions}

We first include the local classification needed for the reduction.

\begin{lemma}\label{lem:coordinates}
Normalize a nonconstant autonomous function on $W_r$ by $f(o)=0$.
Then $f^{-1}(1)$ contains exactly $q$ vertices from each outer triangle,
for one common $q\in\{1,2,3\}$. Conversely every such function is
autonomous, with $a=q/3$ and $b=(4-q)/3$.
\end{lemma}

\begin{proof}
The hub has positive flip probability. Thus no triangle can have all
three values zero, since an outer zero vertex in that triangle would
have zero flip probability. If triangle $i$ has $q_i$ ones, each of its
one vertices has $(4-q_i)/3$ as its flip probability. Constancy on the
one class forces all $q_i$ to agree. Conversely, for a common $q$, both
the hub and every outer zero have flip probability $q/3$, while every
one has flip probability $(4-q)/3$.
\end{proof}

A \emph{three-block packing on $[d]$} is a family of triples
$\{A_i,B_i,C_i\}$ such that each triple consists of pairwise disjoint
subsets with union $[d]$, and all subsets appearing anywhere in the
family are distinct. One block may be empty, but the empty set can occur
at most once in the entire family.

\begin{proposition}[Packing equivalence]\label{prop:packing}
For $d\geq2$, $W_r$ has a separating autonomous code using at most $d$
coordinates if and only if there is a three-block packing of size $r$
on $[d]$.
\end{proposition}

\begin{proof}
Start with a separating code and complement coordinates so that the
hub has the all-zero word. By Lemma~\ref{lem:coordinates}, each column
selects one, two, or three vertices in every triangle. Delete the columns
selecting three; they are constant on all outer vertices, so deleting
them preserves separation among outer vertices. In each column selecting
two, flip its values on the outer vertices only. This is a common
translation of the outer words and also preserves their distinctness.
Every remaining column now has exactly one one in every triangle.

The supports of the three words of a triangle therefore partition the
remaining coordinate set, and all $3r$ supports are distinct. If fewer
than $d$ coordinates remain, add the missing coordinates to one chosen
block in each triple. Restriction to the old coordinates proves that
the blocks remain distinct. This yields a packing on $[d]$.

Conversely, a packing uses $3r$ distinct subsets of $[d]$. Since
$3r\leq2^d$ and $3$ does not divide $2^d$, at least one subset $T$ is
unused. Assign the zero word to the hub and the incidence words of
$A_i\mathbin{\triangle}T$, $B_i\mathbin{\triangle}T$,
$C_i\mathbin{\triangle}T$ to triangle $i$, where $\triangle$ denotes
symmetric difference. Every word is nonzero and the words are distinct.
For a coordinate outside $T$, exactly one vertex per triangle has value
one; for a coordinate in $T$, exactly two do. Each coordinate is
autonomous by Lemma~\ref{lem:coordinates}.
\end{proof}

The unused subset is essential: simply assigning incidence words and
giving the hub zero would fail whenever a block is empty.

\section{Upper bounds by weights on subset sizes}

\begin{lemma}[A weight bound]\label{lem:weight}
Suppose $w_0,\ldots,w_d\geq0$ satisfy
$w_a+w_b+w_c\geq1$ whenever $a+b+c=d$ with $a,b,c\geq0$.
Every three-block packing on $[d]$ then has size at most
\[
 \left\lfloor\sum_{j=0}^d B_jw_j\right\rfloor.
\]
For any real $t>d/3$, a valid choice is
\[
 w_j=\min\left\{1,\max\left\{0,\frac{t-j}{3t-d}\right\}\right\}.
\]
\end{lemma}

\begin{proof}
Each triple has total weight at least one, and no subset is reused.
Summing over its blocks proves the first assertion. For the displayed
weights, if any one of the three untruncated values is at least one,
the corresponding weight is one. Otherwise truncating negative values
to zero can only increase their sum, whose untruncated value is
$(3t-a-b-c)/(3t-d)=1$.
\end{proof}

For $d=3k$, take $t=k+1$. The nonzero weights are
\[
 w_j=1\ (j\leq k-2),\qquad w_{k-1}=2/3,\qquad w_k=1/3.
\]
For $d=3k+1$, the same choice gives
\[
 w_j=1\ (j\leq k-1),\qquad w_k=1/2.
\]
For $d=3k+2$, take $t=k+2$, giving
\[
 w_j=1\ (j\leq k-2),\quad
 w_{k-1}=3/4,\quad w_k=1/2,\quad w_{k+1}=1/4.
\]
Negative indices are omitted. In the last case $B_{k+1}=2B_k$.
Lemma~\ref{lem:weight} therefore gives exactly the three upper bounds
in~\eqref{eq:capacity}.

\section{Constructions attaining the bounds}

We use the following specialization of an established result.

\begin{lemma}[Prescribed block sizes]\label{lem:realization}
Let $(a_i,b_i,c_i)$ be nonnegative integer triples, each summing to $d$.
They can be realized by a three-block packing on $[d]$ if and only if,
for each $j$, the number of occurrences of $j$ among all the triples is
at most $\binom dj$.
\end{lemma}

\begin{proof}
Necessity counts distinct subsets of each size. For sufficiency, apply
Bryant~\cite[Corollary 3]{Bryant} to the hypergraph of all subsets of
$[d]$, with one prescribed color class for each triple. It supplies
distinct subsets with those size counts and makes every color class
almost regular. Each such class has total degree $d$ on $d$ vertices,
so all its degrees are one. Its three subsets thus form a partition.
\end{proof}

To apply the lemma, it suffices to prescribe block sizes and verify the
binomial capacities. We give these lists explicitly. Empty sums below
are zero. We repeatedly use
\begin{equation}\label{eq:binomial}
 \frac{B_{j-1}}{B_j}=\frac{j}{d-j+1},
 \qquad
 B_\ell\geq B_j\quad\text{if }j\leq\ell\leq d-j.
\end{equation}
In the ranges below, the ratios in the lower tail are at most $1/2$.
In particular, $S_{k-3}\leq B_{k-2}$ for $k\geq2$, and
$S_{k-2}\leq B_{k-1}$ for $k\geq1$.

\subsection{The case \texorpdfstring{$d=3k$}{d=3k}}

For each $0\leq j\leq k-2$, prescribe $B_j$ triples of type
\[
 (j,k+1,2k-1-j).
\]
Also prescribe $B_{k-1}$ triples of type $(k-1,k,k+1)$ and
\[
 z=\left\lfloor\frac{B_k-B_{k-1}}3\right\rfloor
\]
triples of type $(k,k,k)$. Their number is
\[
 S_{k-2}+B_{k-1}+z
 =S_{k-2}+\left\lfloor\frac{2B_{k-1}+B_k}{3}\right\rfloor.
\]

The capacities at ranks at most $k$ hold by construction.
For $k\geq2$, the number of blocks of size $k+1$ is
\[
 L=S_{k-3}+2B_{k-2}+B_{k-1}
 \leq3B_{k-2}+B_{k-1}\leq\tfrac52 B_{k-1}\leq B_{k+1}.
\]
The last inequality follows from
$B_{k+1}/B_{k-1}=2(2k+1)/(k+1)\geq3$.
At any larger rank $\ell$, the demand is a single $B_j$ with
$j\leq\ell\leq d-j$, so~\eqref{eq:binomial} applies.
When $k=1$, the list is just one triple of type $(0,1,2)$ and the
capacities are immediate. Lemma~\ref{lem:realization} gives the packing.

\subsection{The case \texorpdfstring{$d=3k+1$}{d=3k+1}}

For each $0\leq j\leq k-1$, prescribe $B_j$ triples of type
\[
 (j,k+1,2k-j),
\]
and prescribe $\lfloor B_k/2\rfloor$ triples of type $(k,k,k+1)$.
This gives $S_{k-1}+\lfloor B_k/2\rfloor$ triples.

Only rank $k+1$ requires a combined capacity check. Its demand is
\[
 L=S_{k-2}+2B_{k-1}+\lfloor B_k/2\rfloor
 \leq3B_{k-1}+\tfrac12B_k < B_{k+1},
\]
because
\[
 \frac{B_{k-1}}{B_k}=\frac{k}{2k+2},
 \qquad
 \frac{B_{k+1}}{B_k}=\frac{2k+1}{k+1}.
\]
Every larger rank receives a single demand $B_j$ covered
by~\eqref{eq:binomial}. All other capacities are immediate.

\subsection{The case \texorpdfstring{$d=3k+2$}{d=3k+2}}

First suppose $k\geq1$. Put
\[
 A=B_{k-1},\qquad B=B_k,\qquad
 x=\left\lfloor\frac{A+1}{4}\right\rfloor,\qquad
 y=B-\left\lceil\frac A2\right\rceil.
\]
For $0\leq j\leq k-2$, prescribe $B_j$ triples of type
$(j,k+2,2k-j)$. Complete the list with
\[
 \begin{array}{c|c}
 \text{block-size triple}&\text{number of copies}\\ \hline
 (k-1,k+1,k+2)&A\\
 (k,k,k+2)&x\\
 (k,k+1,k+1)&y .
 \end{array}
\]
Here $y\geq0$, and checking $A$ modulo four gives
\[
 A+x+y=B+\left\lfloor\frac{3A}{4}\right\rfloor.
\]
The number of triples is therefore the last expression
in~\eqref{eq:capacity}.

At rank $k$ the demand is $2x+y\leq B$; at rank $k+1$ it is
$A+2y\leq2B=B_{k+1}$. For $k\geq2$, the rank $k+2$ demand satisfies
\[
 L=S_{k-3}+2B_{k-2}+A+x
 \leq3B_{k-2}+\tfrac32 A\leq3A < B_{k+2}.
\]
We used $x\leq A/2$, $B_{k-2}\leq A/2$, and
\[
 \frac{B_{k+2}}{A}
 =\frac{2k+3}{k}\cdot2\cdot\frac{2k+1}{k+2}>4
 \quad(k\geq1).
\]
When $k=1$, the rank $k+2$ demand is just $A+x=1$, so it also fits.
Larger ranks again satisfy~\eqref{eq:binomial}.
For $k=0$ (that is, $d=2$), one triple of type $(0,1,1)$ realizes the
claimed capacity one.

All prescribed lists thus meet every binomial capacity. Together with
Lemma~\ref{lem:realization}, the upper bounds of Section~3 and
Proposition~\ref{prop:packing}, this completes the proof of
Theorem~\ref{thm:main}.

\section{A second-order law}

The exact formula identifies the term that was hidden inside the
$O(\log\log r)$ error in the earlier note.

\begin{corollary}\label{cor:asymptotic}
Put $h=H_2(1/3)=\log_2 3-2/3$. Along $d\equiv s\pmod3$,
\[
 M_d\sim c_s\,\frac{2^{hd}}{\sqrt d},
 \qquad
 c_0=\frac7{4\sqrt\pi},\quad
 c_1=\frac9{4\,2^{1/3}\sqrt\pi},\quad
 c_2=\frac{45}{16\,2^{2/3}\sqrt\pi}.
\]
Consequently, as $r\to\infty$,
\begin{equation}\label{eq:second}
 \boxed{\adim(W_r)=
 \frac{\log_2(3r)+\tfrac12\log_2\log_2(3r)}{h}+O(1).}
\end{equation}
\end{corollary}

\begin{proof}
Write $d=3k+s$ with $s\in\{0,1,2\}$. For each fixed $\ell$,
$B_{k-\ell}/B_k\to2^{-\ell}$. The lower-tail ratios are bounded above
by $1/2$, so summing the resulting geometric series gives
$S_{k-2}/B_k\to1/2$ and $S_{k-1}/B_k\to1$.
The three expressions for $M_d/B_k$ therefore tend to
$7/6$, $3/2$, and $15/8$, respectively. Stirling's formula yields
\[
 B_k\sim \frac3{2\sqrt\pi}\,2^{-s/3}\,
              \frac{2^{hd}}{\sqrt d},
\]
which proves the capacity asymptotics.

Let $d=\adim(W_r)$. For large $r$, minimality implies
$M_{d-1}<r\leq M_d$. The three positive constants above are bounded
away from zero and infinity, hence
$\log_2 r=hd-\tfrac12\log_2 d+O(1)$.
This first implies $d\sim(\log_2 r)/h$, and then substituting that
equivalence into the logarithmic term proves~\eqref{eq:second}.
Replacing $\log_2 r$ by $\log_2(3r)$ changes only the bounded term.
\end{proof}

In particular, the logarithmic correction is necessary; it cannot be
replaced by a bounded error in the first-order formula.

\section{Worked thresholds and reproducible certificates}

For seven coordinates, the upper-bound weights are one at subset sizes
zero and one, one half at size two, and zero at larger sizes. Thus
\[
 M_7\leq
 \left\lfloor1+7+\tfrac12\binom72\right\rfloor=18.
\]
The construction uses one triple of sizes $(0,3,4)$, seven of sizes
$(1,3,3)$, and ten of sizes $(2,2,3)$. These require one empty set,
seven singletons, twenty pairs, twenty-four triples, and one four-set.
All demands fit the binomial capacities, so $M_7=18$.

For eight coordinates, the construction uses one triple of sizes
$(0,4,4)$, eight of sizes $(1,3,4)$, two of sizes $(2,2,4)$, and
twenty-four of sizes $(2,3,3)$. Its 35 triples meet the weighted upper
bound. Consequently a windmill with 19 outer triangles, and hence
58 vertices, requires exactly eight autonomous coordinates.

The companion program \texttt{verify\_windmill\_exact.py} performs two
separate checks. First it verifies all prescribed rank demands, and the
matching rational upper-bound certificates, for $2\leq d\leq300$.
Second it constructs actual packings attaining $M_d$ for
$2\leq d\leq10$, converts them to graph codes, and checks distinctness
and all cross-neighbor counts using integer arithmetic.

The construction implements the balancing operation underlying
Bryant's theorem: equal-sized subsets exchange color classes according
to an Euler orientation of an auxiliary multigraph. The code checks
that the sum of squared coordinate degrees strictly decreases at each
step. The resulting partitions, translated vertex codes, and arithmetic
certificates are saved in \texttt{windmill\_exact\_results.json}.
For $d=10$ this yields a code on $W_{116}$, which has 349 vertices.
The $d=3$ witness attains the capacity one but is not a minimum-length
code for $W_1$, whose dimension is two.

These computations are finite checks of the formulas and explicit
witnesses. The proof for all $d$ is the upper-bound argument and the
three capacity-verified constructions above.

\begin{thebibliography}{9}
\bibitem{previous}
\emph{Autonomy Dimension of Graphs: Spectral Separation, Product
Activation, and Entropy Bounds}, research note, October 7, 2026,
Section~7.
\href{https://github.com/geneweng/math-open/blob/dae6532/autonomy_dimension.tex}
{Source at repository commit dae6532}.

\bibitem{Bryant}
D.~Bryant, \emph{On Almost-Regular Edge Colourings of Hypergraphs},
The Electronic Journal of Combinatorics \textbf{23}(4) (2016), Paper~4.7.
\href{https://doi.org/10.37236/5826}{doi:10.37236/5826}.
\end{thebibliography}
\end{document}
